\chapter{Implementation Details}

This chapter describes the implementation details of the Agentic AI-driven Digital Twin framework. It explains how the digital twin, routing execution engine, inventory manager, multi-agent orchestrator, FastAPI server, and React dashboard are implemented in code. It provides file references and outlines key implementation strategies such as multi-key Gemini API rotation, custom agent tools, and dynamic map rendering.

\section{Implementation Overview and File Structure}
The system is built as a decoupled, multi-tier software system using Python 3.11 for all analytical and agentic workloads, and Node.js with React 19 for the frontend presentation. The primary software components are:
\begin{enumerate}
    \item \textbf{Network Model Generator:} Serializes the 14-node, 19-edge directed graph to JSON.
    \item \textbf{Execution Engine and Inventory Layer:} Calculates Dijkstra routing and runs inventory calculations.
    \item \textbf{Agent Orchestration System:} Defines CrewAI agents and tools, using Gemini.
    \item \textbf{FastAPI Backend API Server:} Serves as the REST API endpoint and rotates API keys.
    \item \textbf{React Web Application:} Provides the user dashboard.
\end{enumerate}

\section{Digital Twin Graph Implementation}
The baseline network topology is generated by the script \texttt{network\_model/build\_network\_model.py}. This script initializes a NetworkX \texttt{DiGraph} object and adds 14 nodes and 19 edges with their metadata attributes.

The nodes are defined with attributes such as \texttt{id}, \texttt{label}, \texttt{type}, \texttt{tier}, \texttt{region}, \texttt{role}, and \texttt{disruption\_sensitivity}. Edges represent lanes and store transport mode, transit times, base costs, and friction penalties. The script validates that:
\begin{itemize}
    \item Every node is connected and reachable.
    \item All edge weights are positive.
    \item Alternative routes (e.g., Cape of Good Hope bypass) are modeled.
\end{itemize}
Once validated, the graph is serialized to \texttt{network\_model/outputs/network\_model.json} using a custom JSON exporter that preserves graph properties.

\section{Routing Engine and Shock Injector}
The script \texttt{engine/execution_engine.py} contains the mathematical pathfinding and shock propagation logic. 

\subsection{Shock Propagation Code}
The function \texttt{apply\_shock(graph, event, rules)} takes the baseline graph and applies a shock. It reads rules from \texttt{semantic\_mapping\_rules.json} and computes effective edge transit times and costs. If a node is directly shocked, its edges are scaled by the \texttt{severity\_multiplier}. If a node is a downstream neighbor, the shock is applied with a decayed factor:
\begin{lstlisting}[language=Python]
# Shock propagation logic snippet
if src in direct_nodes or tgt in direct_nodes:
    shock_m = multiplier
elif src in cascade_nodes or tgt in cascade_nodes:
    shock_m = cascade_m # multiplier decayed by factor
else:
    shock_m = 1.0

attrs["effective_transit_time"] = round(base_time * (1 + fp) * shock_m, 4)
attrs["effective_cost"] = round(base_cost * (1 + fp) * shock_m, 4)
\end{lstlisting}

\subsection{Dijkstra Implementation}
Path optimization is executed using \texttt{nx.dijkstra\_path} from the NetworkX library. The destination market is locked during optimization. This ensures that the algorithm only evaluates alternative corridors and supplier origins rather than routing to a different market. The function \texttt{dijkstra\_best\_path} returns the optimal sequence of nodes, total transit time, and total cost.

\section{Stateful Inventory Manager Implementation}
The script \texttt{engine/inventory_manager.py} implements the stateful inventory tracking. 

\subsection{Runtime State Representation}
The inventory manager represents each node's state using a dataclass \texttt{NodeInventoryState}. This tracking is active only for warehouses and SME markets. The class tracks:
\begin{itemize}
    \item \texttt{current\_inventory\_mt}: The current volume of goods in Metric Tons.
    \item \texttt{safety\_stock\_mt}: The volume below which a node is in stockout.
    \item \texttt{daily\_demand\_mt}: The daily rate of consumption (non-zero only at markets).
    \item \texttt{reorder\_point\_mt} and \texttt{reorder\_quantity\_mt}: Parameters for the $(s, Q)$ ordering policy.
    \item \texttt{pending\_orders\_mt}: Inbound replenishment orders that are in transit.
\end{itemize}

\subsection{Demand and Order Propagation}
A simulation step is advanced daily via the \texttt{advance\_day()} method. When a warehouse's inventory falls below its reorder point, it places an order upstream. To model the Bullwhip Effect, warehouses amplify their reorder quantity by a factor proportional to safety stock:
\begin{equation}
    Q_{order} = Q_{reorder} \times \left(1 + \frac{I_{safety}}{Q_{reorder}}\right)
\end{equation}
The inventory manager tracks order history and calculates the bullwhip ratio (variance of orders divided by variance of demand) for each warehouse.

\section{CrewAI Agent Orchestration and Custom Tools}
The multi-agent orchestration is implemented in \texttt{agent\_orchestration/main.py}. The four agents are defined as CrewAI \texttt{Agent} objects using Google's Gemini Flash model (\texttt{gemini/gemini-flash-latest}).

\subsection{Custom Agent Tools}
Three custom tools are implemented using the \texttt{@tool} decorator to allow the Cost Analyst agent to inspect and interact with the stateful inventory system:
\begin{enumerate}
    \item \textbf{CheckInventoryTool:} Takes a \texttt{node\_id} and returns a JSON snapshot of the node's current inventory, days of cover, and stockout status.
    \item \textbf{NetworkInventorySnapshotTool:} Returns a network-wide inventory health summary (days of cover, utilization, stockout status) for all active nodes.
    \item \textbf{EmergencyRestockTool:} Restocks a node by adding inventory up to its maximum capacity. The analyst uses this tool when a delay will cause a stockout, modeling emergency air-freight logistics.
\end{enumerate}

\subsection{Sequential Workflow Assembly}
The agents and tasks are assembled into a CrewAI \texttt{Crew} executing with \texttt{process=Process.sequential}. A file-logger saves the intermediate decisions to \texttt{orchestration.log} and the execution details to \texttt{crewai\_execution.log}.

\section{FastAPI REST Backend API Server}
The FastAPI backend server is implemented in \texttt{logistics\_app_backend/main.py}.

\subsection{API Key Rotation Middleware}
To prevent hitting Gemini API rate limits, a round-robin key rotation middleware is implemented. The server reads all environment variables starting with \texttt{GEMINI\_API\_KEY\_} and rotates to the next key before executing any agent crew:
\begin{lstlisting}[language=Python]
_ALL_KEYS = [
    v for k, v in sorted(os.environ.items())
    if k.startswith("GEMINI_API_KEY_") and v.strip()
]

def _next_key() -> str:
    global _key_index
    with _key_lock:
        key = _ALL_KEYS[_key_index % len(_ALL_KEYS)]
        _key_index += 1
        return key
\end{lstlisting}
If a request encounters a \texttt{RESOURCE\_EXHAUSTED} (429) rate limit error, the middleware catches the error and retries the request with the next key.

\subsection{API Endpoints}
The server exposes two primary REST endpoints:
\begin{itemize}
    \item \texttt{GET /api/shock}: Calls the Friction Sentry agent to generate a fresh, dynamic maritime disruption scenario.
    \item \texttt{POST /api/optimize}: Receives an optimization request containing \texttt{optimize\_by} (time or cost). It triggers the CrewAI workflow, optimizes routes, and returns a JSON payload containing the optimal route node IDs, deltas, and visual flags.
\end{itemize}

\section{React Dashboard UI Implementation}
The presentation layer is implemented as a React web application in \texttt{logistics\_app\_frontend/}. 

\subsection{Interactive Map Component}
The dashboard uses Leaflet (\texttt{react-leaflet}) to render the supply chain network on a geographical map. The 14 nodes are plotted at their physical coordinates. The active route is drawn as a colored polyline. When a shock occurs, the affected corridor is highlighted in red, the baseline path in green, and the optimized path in blue.

\subsection{UI Components and State Management}
Axios is used to make API calls to the FastAPI backend. Framer Motion handles smooth transitions and alerts. Lucide React provides system icons. The main component \texttt{App.jsx} manages the UI state, displaying:
\begin{itemize}
    \item Cost and time delta widgets.
    \item Inventory health trackers.
    \item A timeline of shock propagation.
    \item Raw agent dialogue transcriptions.
\end{itemize}

\section{Summary}
This chapter detailed the code implementations across all five layers, including the graph model scripts, Dijkstra solvers, stateful inventory tracking, multi-agent frameworks, FastAPI backend server with key rotation, and React frontend dashboard. The next chapter presents the verification testing and experimental results of this implementation.
